\section*{Carl Friedrich Gauss: Werke IX}
\begingroup\footnotesize
\noindent\textit{Werke}, Band IX, Nachlass. \EditionScope.

The reader contains the historical outer cover, half-title and title with
its portrait. The determination of the difference of latitude between the
observatories of Göttingen and Altona includes its title and contents,
introduction, 43-star catalogue and observation series. The Resultate include
the weighting formulas, collimation errors, residuals, statistical reductions,
direct and reflected observations, and star comparisons. The Seeberg section includes its sixty
observation results, followed by the Zusatz with twenty-five survey
localities and the Anzeige with its historical remarks.
Erdellipsoid
und geodätische Linie includes its title and blank reverse, the three
Nachlass articles on the ellipsoid with their numerical calculations
and remarks, and Articles [1]--[13] on the geodesic line. These include
the auxiliary tables, the two 61-row tables, the worked calculations,
the editable diagrams and the historical remarks on pages 96--98.
The letter to J. J. Baeyer on the variation of latitude with height occupies
pages 99--102, with its editable diagram and historical remarks.
The following two
Nachlass notes concern spherical triangle angles and points on a great
circle, with an editable diagram and remarks signed Krüger and Börsch.
Conforme Doppelprojection includes its title and blank reverse, Articles
[1]--[10] on transferring the ellipsoidal surface to the sphere, the
numerical table and historical remarks, followed by stereographic
projection of the sphere onto the plane and its remarks.
The Mercator-projection component
contains Articles [1]--[8], the numerical table and the historical remarks
signed Krüger and Börsch. The stereographic representation of the spheroid
contains the formulas and historical remarks on pages 133--134. The conical
mapping component contains its title and blank reverse, the Nachlass with
its formulas and two numerical tables, and the historical remarks.
The conformal-mapping component contains its title and blank reverse, the opening Nachlass heading,
Articles [1]--[32] and the historical Bemerkungen.
Its coordinate transformations, logarithmic and curvature formulas,
numerical tables and worked Varel calculation are represented as editable text.
The Landesvermessung component contains the letters to Schumacher of
18 April 1830, 30 April 1830, 17 May 1831, 25 June 1831 and 9 December 1838,
including the auxiliary tables and Neuwerk calculations, followed by the
historical remarks signed Krüger and Börsch. The trigonometrical
point-determination component contains its title and blank reverse,
Articles [1]--[9] and historical remarks, with editable diagrams, tables,
the Goldschmidt quotation and Krüger credit.
Ausgleichung einfacher Figuren includes its title and blank reverse,
Articles [1]--[7], the letters of 11 February 1824 and 19 January 1840,
and the remarks signed Krüger, with editable diagrams, formulas and tables.
The Stationsausgleichungen component includes its title and blank reverse,
the Zeven, Brillit, Wilsede and Windberg observations and normal equations,
with the historical Hartmann credit, the Breithorn observations of 13--16 September 1822, and the letters
to Gerling of 26 December 1823 and to Schumacher of 22 December 1827
and 7 January 1828. Their tables, triangle drawing and historical notes
are followed by the Bemerkungen signed Krüger, including the quotation
from Goldschmidt's Geodätischer Calcül nach Gauss. The opening of Zur Netzausgleichung
contains its title and blank leaves, the condition-count formulas [1],
the 33 connection-list entries and editable network with 75 edges [2],
and all 51 observed-triangle groups and concluding relations [3].
This component also contains Article [4], including its 51 numbered equations,
the opening and all twelve worked cases I--XII of Article [5], Article [6] with its
relation table, correction expressions and twelve normal equations,
Article [7] with its 150 direction corrections and concluding text,
Article [8] with its 51 triangle groups and Hamburg--Hohenhorn connection table,
Article [9] with its 150 azimuth and reduction rows,
and the six letter excerpts in Article [10]: Gauss to Olbers,
2 November 1823, 6 July 1824 and 14 May 1826, and Gauss to Gerling,
2 May 1837, 5 June 1838 and 14 November 1838, including the historical footnotes,
the six-point triangulation diagram and the quadrilateral ring diagram,
and the complete historical Bemerkungen, signed Krüger.
The Oldenburg component contains its title leaf and blank reverse, opening paragraph,
84 preliminary table rows in 21 station groups, and the triangulation
diagram with 21 numbered vertices and 42 edges, the 22 angle-sum errors,
the condition examples, all 22 normal equations and correlate values,
the correction discussion, the complete triangle overview with 63 angle
and side-logarithm rows, and the coordinate formulas and table with
23 coordinate rows and two difference rows, and the second-adjustment
formulas with six direction relations, the 77 coefficient pairs and
their explanation, and the complete historical Oldenburg remarks,
signed Krüger, including the quotation from Gauss's letter to v. Schrenck
of 3 August 1835. The historical adjustment itself is unfinished.
The Hannover component contains its title leaf, blank reverse and the eight letter
excerpts from Gauss to Schumacher in [1], 5 July 1816 to 14 January 1829,
including all three historical footnotes, and all eight dated excerpts
in [2], Gauss to Bessel, 26 December 1821, 15 November 1822,
5 November 1823, 20 November 1824, 12 March 1826, 20 November 1826, 1 April 1827 and 9 April 1830, including their three historical footnotes.
The Hannover component also contains the letter from Gauss to Bohnenberger in [3],
16 November 1823, including its three historical footnotes, and all twelve
Gauss-to-Olbers excerpts in [4], dated 13 January 1821, 18 April 1822, 4 July 1824, 8 July 1824, 19 February 1825, 25 February 1825, 9 October 1825, 2 April 1826, 14 January 1827, 1 March 1827, 14 June 1830 and 2 September 1837,
including three historical footnotes, the triangulation diagram labelled Bremen, W., Brüttendorf and Bottel,
and the 57-row \textit{Relative Höhen} table.
All eleven Gauss-to-Gerling excerpts in [5] are included: 5 October 1821,
21 February 1822, 7 November 1822, 27 July 1823, 11 August 1823, 1 September 1823, 5 September 1823, 3 October 1823, 19 July 1827, 12 September 1838 and 14 November 1838,
including the four heliotrope signalling diagrams, three optical diagrams,
coordinate table and historical footnote, followed by the Bemerkungen signed Krüger.
Section [II.] contains the two published reports of 1821 and 1822,
Schumacher's introduction, the coordinate table and the historical Bemerkungen.
Section [III.1] contains the plan of the work, its introduction and the
opening of its first section. Section [III.2] contains the excerpts from
the reports of 7 January 1822, 31 January 1823, 16 February 1825,
21 November 1827 and 26 June 1828, the historical report dated
8 February 1838, and the excerpts of 5 July 1840 and December 1844,
followed by the complete thirty-three-row coordinate table [III.3]
and the historical Bemerkungen signed Krüger. The associated foldout
map, devised by Lieutenant J. Gauss, is represented as an editable
vector drawing with its triangulation networks, geographical outlines,
legends and scale. The accompanying historical note identifies
Krüger's editorial corrections to the map.
The Höhenmessungen component contains its title leaf and blank reverse,
the three refraction texts [1]--[3], including the letters to Schumacher
of 28 November 1824 and 27 December 1846, and the historical editorial note.
It also contains the Nachlass on terrestrial refraction and height adjustment,
including the 71 error equations, 31 normal equations, station corrections
and historical Bemerkungen signed Krüger, and the three tables with
instructions for barometric height measurement.
The Heliotrop component contains its title leaf and blank reverse,
the Göttingische gelehrte Anzeigen article of 9 August 1821 with its optical
diagram, the Astronomisches Jahrbuch contribution submitted on 26 December
1821, the letter to Olbers dated 1 July 1821 and four letters to Schumacher
of 11 July 1821, 8 November 1821, 10 July 1822 and 15 January 1827,
including their notes and drawing. It contains the article on heliotrope
adjustment, its thirteen-figure plate, the two Nachlass notes with their
formulas and geometric drawing, and the historical Bemerkungen signed Krüger
with the diagram of the first construction.
The Messungsfehler component contains its title leaf and blank reverse,
the text on surveying instruments, the letters on measurement errors
to Olbers (July 1825), Schumacher (14 August 1825), Bessel
(29 October 1843 and 15 August 1844) and Gerling (17 April 1844),
the observation calculations and historical Bemerkungen signed Krüger.
The following component contains its title leaf and blank reverse,
the three notes on reducing inclined angles to the horizon, their worked
example, calculation scheme and historical Bemerkungen, and
Repetitionsbeobachtungen, including the letter to Bessel of 27 January
1819, the recurrence and error formulas and the historical remarks signed
Krüger. The Bemerkungen zum neunten Bande signed F. Klein, historical
contents, final blank leaf and publisher advertisement are included.
The historical contents retain the original page references.

\par\endgroup
\clearpage
\begingroup\small
Source: the historical edition digitized by the Göttinger
Digitalisierungszentrum.\par
\noindent Identifier PPN23601515X; printed pages
1--528, including the unnumbered Hannover foldout map and the Conforme Doppelprojection des Sphäroids auf die Kugel und die Ebene, Conforme Übertragung des Sphäroids auf den Kegelmantel, Conforme Abbildung des Sphäroids in der Ebene, Trigonometrische Punktbestimmung, Ausgleichung einfacher Figuren, Stationsausgleichungen, Netzausgleichung, Oldenburg, Hannover, Höhenmessungen, Heliotrop, Messungsfehler and Reduction schiefer Winkel title leaves and the blank leaves at pages 66, 106, 136, 142, 220, 244, 262, 294, 296, 330, 344, 436, 460, 486, 506 and 508, together with the final unnumbered blank leaf and publisher advertisement. The historical editorial wording, notes,
numerical values and signs are retained.

\noindent\href{https://gdz.sub.uni-goettingen.de/download/pdf/PPN23601515X/PPN23601515X.pdf}{Historical source scan, Göttinger Digitalisierungszentrum}.

Comparison scan: Google Books digitization, Internet Archive identifier:\\
\texttt{carlfriedrichga01gausgoog9}.

The separate point before den on page 51 follows the principal copy;
its function is uncertain. The gathering signature on that page follows
the Google copy. The compact cross-shaped Greek designation in star row
8 is represented by the kappa variant, with its precise historical
typeform qualified. Numerical and geographical spellings in the Seeberg
and survey tables follow the printed text, including unreduced seconds.

The historical title-page portrait is reproduced from the principal scan
as an original illustration; its creator is not named on the inspected leaf.
It is not AI-created. All title text, body text, mathematics, tables and
mathematical diagrams are editable LaTeX. Library ownership stamps,
accession handwriting and the modern digitizer cover are copy-specific
additions outside the historical edition and are not reproduced.

The two low marks in the running head on page 6 and the trailing mark
after 54,24 in the page 18 observation table follow the principal copy;
their functions are uncertain. The star designation on page 9 uses the
same qualified kappa variant as page 51. The Altona colon on page 12 and
Jun without a period on page 15 follow the principal text.

The principal copy's separate point after beta-prime in the weighting
formula on page 30 is retained; its function is uncertain. The differing
latitude values on pages 38 and 39 follow the printed text without correction.
The signature on page 35 is supplied from the Google and HKUST copies.

Faint citation delimiters on pages 64 and 71, small digits on page 69,
formula delimiters and the outer square on page 70, and a faint
parenthesis on page 72 follow the clearer HKUST copy. The separate point
before the final logarithm on page 69 follows the principal copy; its
function remains uncertain. The auxiliary table on page 77 preserves
the printed blank cells and abbreviated values. Diagrams on pages 73
and 75 preserve their labelled connections with approximate geometry.

The clipped rightmost column of the table on page 82 follows the HKUST
copy. The flattening values on pages 82 and 84 differ from those on
pages 83 and 87, as printed. The separate points on pages 84, 88, 93 and
94 and the raised point in the formula ending on page 91 are retained
without assigning them an additional mathematical function. A faint
closing parenthesis on page 86 and the fourth power on page 98 follow
the clearer HKUST copy. The diagrams on pages 100 and 103 use approximate
geometry while preserving their labelled connections and crossings.

The page range 30--32 cited on page 102 and small notation on pages
109--112 follow the clearer HKUST copy. The separate point before the
denominator on page 113 follows the principal copy; its function is
uncertain. The vertical stroke after cos on page 119 is present in the
principal, Google and HKUST copies; its function remains uncertain.
Unclassified peripheral ink on pages 112, 117--118 and 120--121 is not
represented as additional mathematical characters. The page 103 diagram
preserves the printed connections and labels with approximate geometry.

The lowered point beneath 2345.09 on page 126 is retained separately
from the digits; its function is uncertain. The short stroke near 8504.24
is unclassified. A faint closing parenthesis on page 129 follows the
clearer HKUST copy. That copy also clarifies the denominator 45360 and
delimiters on page 132 and the series on page 138.
On page 139 the gaps in 6,709 8436 473 and 1180 553 follow the principal
copy; Google has points at those positions. The separate point before
860 in the lower table follows the principal and is absent from Google.
The separate points on pages 138--140 are retained with uncertain functions.
Unclassified peripheral ink on pages 132, 137 and 140 is not represented
as additional mathematical characters.

Faint closing parentheses on page 144 follow the clearer HKUST copy;
those in the first log hyp D formula on page 147 follow Google.
On page 160 the point after 89,437 follows the principal copy and is absent
from Google and HKUST. The two marks before 109 on page 161 and the
raised quote-like mark before 89 on page 162 are retained with uncertain
functions. The table value 788 on page 162 and denominator 5040 on
page 171 follow the clearer comparison witnesses. The separate points
on pages 163, 167, 169--171 follow the principal copy; their functions
are uncertain. Unclassified peripheral ink on pages 144, 147--150 and
166 is not represented as additional mathematical characters.

On page 195 the comma after zugefügt worden follows the principal copy;
Google and HKUST have a semicolon. On page 197 the period after the first
displayed 3,418 9411.803 follows the principal; the other copies have a comma.
The small coefficient on page 195 and cramped notation on pages 196--199
and 201--202 are supported by the clearer HKUST copy. The unclassified
isolated marks on pages 192 and 198 and in the otherwise blank degree cell
of the page 210 table are not represented as additional characters.
The separate points in the calculation field on page 203, after Es and
beside the first longitude equation on page 208, and above the minus
in formula 6) on page 209 follow the principal copy; their functions are uncertain.

On page 212 the exponent 2 in the denominator of F follows the clearer
HKUST copy. The upper and lower marks between 434 and 2945 in k
follow the principal copy; their functions remain uncertain. The separate
short stroke after the comma following gebrauchen on page 225 follows
the principal copy and is absent from Google; its textual status is uncertain.

On page 233 the principal semicolon after scheint differs from the counterpart comma;
the raised comma-like mark by the title on page 239 has uncertain function.
On page 242 the principal period after sein differs from the counterpart colon.
Its point before Rechnungen and the raised point before eine on page 249 are
retained with uncertain function; obscured Google page 250 text follows the principal.
The principal scan clips the right-hand part of the system on page 270;
the complete Google copy supplies that portion and the sideways orientation.
Unreadable areas of Google page 254 follow the principal copy. Separate marks
on pages 254--257, 260, 263--264 and 268--269 follow the principal; their
functions remain uncertain. The repeated e headings on page 265 are retained.

The separate marks beside the formulas on pages 272--273, the folios
on pages 274 and 276, the value 253 on page 276, the footnote on
page 278 and the two prose/display positions on page 281 follow the
principal copy; their functions and physical origins are uncertain.
The closing rule on page 277 in Google is absent from the principal
copy and is not imported.
On page 286 the first item has a comma after I in the principal copy;
the comparison copy has a point. The inferior logarithmic marks on
page 287 are represented by n; their exact printed type is uncertain.
The separate marks on pages 292 and 297 and in the Wilsede row on
page 298 follow the principal copy; their functions are uncertain.
The first Wilsede leader is comma-shaped in the principal copy and
point-shaped in HKUST. On page 300 the triangle 20 sum has a point
in 0.545, whereas Google has a comma; the separate point in triangle 8
and the dash below the middle table follow the principal positions.
The upper point above the comma in triangle 48 on page 301 follows
the principal copy. Counterpart-only isolated marks on pages 283, 293
and 301 are not imported. Unidentified peripheral ink on pages 282--285
is not represented as additional characters. The page 299 network
preserves the printed incidences and labels with approximate geometry.

An additional copy at HKUST Library supports the clipped table on page 82 and small notation on pages 86, 98, 102, 109--112, 119, 125, 129, 132, 138, 144,
152--162, 171 and 195--199,
201--202 and 212, the network and connection list on
pages 298--299 and both \(\tfrac13\,\text{Excess}\) table headings on pages 304--305 and the damaged letter in
\textit{20-maliger} on page 441.\par
\noindent\href{https://lbezone.hkust.edu.hk/rse/?p=13258}{Comparison copy, HKUST Library}.\par
\noindent\href{https://lbezone.hkust.edu.hk/obj/4/o/b761099/v9/pg00220.jpg}{Page 212, HKUST Library}.
The signature on page 465 and the terminal rule on page 466 follow the
Google comparison copy; they are absent from the principal scan.
The isolated points on pages 473, 474 and 480 follow the principal copy;
their function is uncertain. The plate labels on page 477 retain the
printed J and lower-case c, differing from the prose references I and CC.
The semicolon after Müller on page 489 and the raised point on page 495
follow the principal copy; Google has a comma and no raised point, respectively.
The terminal point after Bremervörde in the page 492 diagram is uncertain
in the principal copy and dot-like in Google. The irregular ink attached
to the readable word Ich on page 495 is not reproduced as an additional sign.
The separate points before Fernrohr on page 501 and Fehler on page 515
follow the principal copy and are absent from Google; their function is uncertain.
The signature IX. / 64 on page 505 follows Google and is absent from the
principal scan. Faint parentheses on page 512 follow the readable comparison
copy. The bracket-like mark between the equals sign and cube root in the
definition of mu on page 513 is represented by ]; its function is uncertain.
Irregular peripheral ink near Aus and the final exponent on page 516 is
not represented as additional text.
On page 519 the exponent 2 in the elimination equation and exponent 4
in the reciprocal-weight numerator follow the principal copy; the Google
superscripts are filled. The date 18 September 1822 on page 526 follows
the principal copy; Google appears to read 19. The separate point after
Auszug on page 526 is retained. Unidentified peripheral ink near the
heading on page 521, Azimuth on page 524 and an Olbers on page 526 is
not represented as additional characters. On page 528 the unaccented Uber
in Uber Messungsfehler and the absent colons after Messungsfehler and
Repetitionsbeobachtungen follow the principal copy. The extra point in
the 1822 Juli 10 leader and comma-shaped mark before heading B. in the
advertisement are retained with uncertain functions. The advertisement
follows the principal final leaf; it has no counterpart in the Google copy.

The unfolded Hannover map follows the separate digital copy at
SUB Göttingen; the principal PDF omits it, and the Google PDF shows
folded strips only. Its geometry, stroke weights and colour-band offsets
are approximations to the printed drawing. Overprinted contour
centre-lines at Pugelatz, Fuhrenberg and the Inselsberg ray cannot be
distinguished throughout; no hidden continuation is supplied.\par
\noindent\href{https://webdoc.sub.gwdg.de/ebook/e/2005/gausscd/html/kapitel_landv_karte.htm}{Hannover triangulation map, SUB Göttingen}.

The quoted letter also appears in the \textit{General-Bericht über die
mitteleuropäische Gradmessung für das Jahr 1865}, Berlin: Georg Reimer,
1866, pp.~23--24.\par
\noindent\href{https://books.google.com/books?id=VBNDAAAAcAAJ\&pg=PA24}{Earlier printed letter, Google Books}.

The editable text and mathematical typesetting are AI-transcribed
and AI-produced. Modern fonts, spacing and line breaks represent
the historical content without reproducing the original type.

\par\endgroup
